\documentclass[11pt]{amsart}

\usepackage{amsmath}
\usepackage{amssymb}
\usepackage{amsthm}
\usepackage{booktabs}
\usepackage{array}
\usepackage{hyperref}

\newtheorem{theorem}{Theorem}[section]
\newtheorem{proposition}[theorem]{Proposition}
\newtheorem{lemma}[theorem]{Lemma}
\newtheorem{corollary}[theorem]{Corollary}
\theoremstyle{definition}
\newtheorem{definition}[theorem]{Definition}
\theoremstyle{remark}
\newtheorem{remark}[theorem]{Remark}

\numberwithin{equation}{section}
\setlength{\emergencystretch}{2em}

\newcommand{\NN}{\mathbb{N}}
\newcommand{\ZZ}{\mathbb{Z}}
\newcommand{\QQ}{\mathbb{Q}}
\newcommand{\RR}{\mathbb{R}}
\newcommand{\CC}{\mathbb{C}}
\newcommand{\cO}{\mathcal{O}}
\newcommand{\cD}{\mathcal{D}}
\newcommand{\cC}{\mathcal{C}}
\newcommand{\cP}{\mathcal{P}}
\newcommand{\Nm}{\operatorname{N}}
\newcommand{\eu}{\phi_{\mathrm E}}
\newcommand{\Res}{\operatorname{res}}
\newcommand{\lean}[1]{{\def\_{\textunderscore\allowbreak}\texttt{#1}}}

\begin{document}

\hypersetup{pdftitle={The reciprocal cube sums of Lucas sequences of quadratic units are irrational},pdfauthor={Shawn Krisman}}
\title[Reciprocal cubes of Lucas sequences]{The reciprocal cube sums of Lucas sequences of quadratic units are
irrational}

\author{Shawn Krisman}

\subjclass[2020]{Primary 11J72; Secondary 11B39, 33D15, 41A21, 68V20}
\keywords{Lucas sequences, Pell numbers, Fibonacci numbers, irrationality, $q$-series, Hankel determinants, Lean}

\begin{abstract}
Let $\eta>1$ be a unit of a real quadratic order and $U_n=(\eta^n-\bar\eta^n)/(\eta-\bar\eta)$ its Lucas sequence. We
prove that $\sum_{n\ge1}U_n^{-3}$ is irrational. Equivalently, $\sum_{n\ge1}U_n(P,Q)^{-3}\notin\QQ$ for the Lucas
sequences with $Q=-1$, $P\ge1$, and with $Q=1$, $P\ge3$. In particular the sums of the reciprocal cubes of the Pell
numbers and of the numbers $F_{kn}$, for every $k\ge1$, are irrational. The case of the Fibonacci numbers was treated in
earlier work of the author; the proof here is the same Hankel-determinant construction, carried out for a general unit
and for both signs of its norm. The constants of the argument do not depend on $\eta$ once they are measured in units
of $\log\eta$. The whole argument is formalised in Lean~4 with Mathlib.
\end{abstract}

\maketitle

% ===========================================================================
\section{Introduction}
% ===========================================================================

For integers $P,Q$ let $U_n=U_n(P,Q)$ be the Lucas sequence $U_0=0$, $U_1=1$, $U_{n+1}=PU_n-QU_{n-1}$. When $Q=\pm1$
and $D=P^2-4Q>0$ is not a square, $U_n=(\eta^n-\bar\eta^n)/\sqrt D$ with $\eta=(P+\sqrt D)/2$ a unit of the order
$\ZZ[\eta]$ of norm $\eta\bar\eta=Q$. Conversely, every unit $\eta>1$ of a real quadratic order arises in this way, with
$P=\eta+\bar\eta\ge1$ if $Q=-1$ and $P\ge3$ if $Q=1$. The Fibonacci numbers are $U_n(1,-1)$, the Pell numbers $\cP_n$ are
$U_n(2,-1)$, and $F_{kn}/F_k=U_n(L_k,(-1)^k)$, where $L_k$ is the $k$-th Lucas number.

\begin{theorem}\label{thm:main}
Let $Q=-1$ and $P\ge1$, or $Q=1$ and $P\ge3$. Then $\displaystyle\sum_{n\ge1}\frac1{U_n(P,Q)^3}$ is irrational.
\end{theorem}

\begin{corollary}\label{cor:main}
The numbers $\sum_{n\ge1}\cP_n^{-3}$ and, for every $k\ge1$, $\sum_{n\ge1}F_{kn}^{-3}$ are irrational.
\end{corollary}

\begin{proof}
Theorem~\ref{thm:main} with $(P,Q)=(2,-1)$, and with $(P,Q)=(L_k,(-1)^k)$, which is admissible: $L_k\ge1$, and $L_k\ge3$
for $k\ge2$. Then $U_n=F_{kn}/F_k$ by Binet's formula, so $\sum_nF_{kn}^{-3}=F_k^{-3}\sum_nU_n^{-3}$.
\end{proof}

The case $(P,Q)=(1,-1)$, $\sum_nF_n^{-3}\notin\QQ$, is the main result of \cite{Kri26}; we refer there for the history
of the Fibonacci zeta function. Briefly: Andr\'e-Jeannin proved that $\sum_nF_n^{-1}$ is irrational \cite{AJ89}, and
Duverney, Ke.~Nishioka, Ku.~Nishioka and Shiokawa proved that the even powers $\sum_nU_n^{-2s}$ are transcendental for
the generalised Fibonacci sequences with $Q=-1$, as are the odd-index sums $\sum_nU_{2n-1}^{-s}$ \cite{DNNS97}; the
algebraic relations among such numbers were found by Elsner, Shimomura and Shiokawa \cite{ESS07,ESS08}. These sums are
values of quasi-modular forms, to which Nesterenko's theorem \cite{Nes96} applies. The odd powers of the full sums are
not of this kind (see \cite{Wal10}). For the Pell numbers, the tails of $\sum_n\cP_n^{-3}$ have been studied by Xu and
Wang \cite{XW13}, who determined the integer part of their reciprocals; we know of no irrationality result for any
sum $\sum_nU_n^{-3}$ before \cite{Kri26}.

\subsection*{The method}
Put $t=\eta^{-1}$ and $w=Qt^2$. Since $\bar\eta/\eta=Qt^2$, we have $U_n=t^{-n}(1-w^n)/\sqrt D$, and so
\begin{equation}\label{eq:Xi}
\Xi:=\frac1{D\sqrt D}\sum_{n\ge1}\frac1{U_n^3}=\sum_{N\ge1}\frac{t^{3N}}{(1-w^N)^3}.
\end{equation}
We follow \cite{Kri26}, which adapted to $\ZZ[\varphi]$ a construction of an unrefereed research draft on the $q$-zeta
value $\sum_nn^2/(p^n-1)$, $p\ge2$ an integer \cite{UPL,QZ3}; the construction is a moving-denominator Hankel
determinant in the line of Uvarov \cite{Uva69}, Krattenthaler \cite{Kra21}, Pr\'evost \cite{Pre96} and Zudilin
\cite{Zud16,Zud17}. A linear functional $L$, a second Euler derivative summed against the weights $(Qt)^n$ at the nodes
$w^n$, has moments $L(s^k)\in\QQ(\sqrt D)$, rational functions of $t$, and satisfies $L(1/(1-w^as))=\lambda_a(\Xi-S_a)$
with explicit $\lambda_a$ and partial sums $S_a$ of \eqref{eq:Xi}. The Hankel determinants $\Delta_K(X)$ of $L$ against
$D_K(s)=\prod_{a\le K}(1-w^as)$ satisfy $\Delta_K(\Xi)=\det[L(s^{i+j}/D_K)]$, and four estimates are combined:
\begin{center}
\small
\begin{tabular}{lll}
\toprule
 & size (per $K^3\log\eta$) & \\
\midrule
real decay $|\Delta_K(\Xi)|$ & $-2/9$ & Proposition~\ref{prop:decay}\\
conjugate $|\sigma\Delta_K(-\Xi)|$ & $-1$ & Proposition~\ref{prop:conj}\\
denominators $\log|\Nm C_K(t)|$ & $\le c_*=1.1954\ldots<11/9$ & Propositions~\ref{prop:global}, \ref{prop:cost}\\
nonvanishing & infinitely many $K$ & Proposition~\ref{prop:nonvan}\\
\bottomrule
\end{tabular}
\end{center}
If $\sum_nU_n^{-3}$ were rational, $C_K(t)(mD^2)^K\Delta_K(\Xi)$, where $m$ is the denominator of the sum, would be a nonzero element of $\ZZ[\eta]$ of norm $<1$
for large $K$.

Compared with \cite{Kri26} there are three changes. First, the sign: for $Q=1$ the nodes $w^n=t^{2n}$ are positive and the
roots of unity at which the coefficients have poles are of a different type (Section~\ref{sec:den}); the local
arguments are unchanged, only the classification of the roots changes, and the elementary count of the $\alpha$-roots
used in the proof is the same for both signs (the difference appears only in the sharp constant of
Remark~\ref{rem:sharp}). Second,
the unit: all analytic estimates use only $0<t\le\varphi^{-1}$, which holds for every admissible $(P,Q)$, and every
constant is a nondecreasing function of $t$ on that interval. Third, the order $\ZZ[\eta]$ need not be maximal; the
argument uses only that the norm of a nonzero element of $\ZZ[\eta]$ is a nonzero rational integer. Sections
\ref{sec:setup}--\ref{sec:proof} follow \cite{Kri26} closely, often verbatim with $\varphi^{-1}$ replaced by $t$ and
$-1$ by $Q$; we repeat the arguments so that the present paper is self-contained.

The same argument gives $\sum_nU_{kn}^{-3}\notin\QQ$ for every admissible $(P,Q)$ and $k\ge1$, since
$U_{kn}=U_kU'_n$ with $U'$ the Lucas sequence of the unit $\eta^k$. The normalisation $\eta>1$ does not cover the
alternating sums $\sum_n(-1)^{n-1}U_n(P,-1)^{-3}=\sum_nU_n(-P,-1)^{-3}$.

\subsection*{Formal verification}
The theorem and every step of its proof are formalised in Lean~4 with Mathlib (Section~\ref{sec:lean}). The analytic
ingredients are formalised for every real $t$ with $0<t\le\varphi^{-1}$ and both signs, and the integrality statement
for every real $\tau\ne0$ that is not a root of unity, with one integer polynomial for all such $\tau$; only the final
assembly uses that $t$ is the inverse of a quadratic unit.

% ===========================================================================
\section{The functional and the determinants}\label{sec:setup}
% ===========================================================================

\subsection*{Notation}
Throughout, $(P,Q)$ is admissible: $Q=-1$ and $P\ge1$, or $Q=1$ and $P\ge3$. Put $D=P^2-4Q$, $\eta=(P+\sqrt D)/2$,
$\bar\eta=(P-\sqrt D)/2=Q/\eta$, $t=\eta^{-1}$, $\cO=\ZZ[\eta]$, and let $\sigma$ be the nontrivial automorphism of
$\QQ(\sqrt D)$, so $\sigma t=1/\bar\eta=Q\eta$. We write $\varphi=(1+\sqrt5)/2$, Euler's function $\eu(d)$ and the
M\"obius function $\mu(d)$. Put
\[
\rho=Qt,\qquad w=Qt^2,\qquad \alpha_k=\rho w^k\ (k\ge0),
\]
\[
\psi_N=\frac{t^{3N}}{(1-w^N)^3},\qquad S_a=\sum_{N=1}^a\psi_N,
\]
$A_K=K(K+1)/2$ and $\theta=s\,\frac{d}{ds}$. Then $\rho^2=Qw$, $|w|=t^2$, $\Xi=\sum_N\psi_N$, and $w/\rho=t$. We write
$h_m$ for the complete homogeneous symmetric polynomial of degree $m$ ($h_m=0$ for $m<0$), and $K\ge2$.

\begin{lemma}\label{lem:unit}
\textup{(i)} $D$ is not a square, so $\sqrt D\notin\QQ$. \textup{(ii)} $\eta\ge\varphi$, so $0<t\le\varphi^{-1}$.
\textup{(iii)} $t$ and $\sigma t$ are the roots of $\tau^2-PQ\tau+Q$, and $t\,\sigma t=Q$.
\textup{(iv)} $\sqrt D=P-2Qt$ and $P-2Q\,\sigma t=-\sqrt D$.
\textup{(v)} $U_n=(\eta^n-\bar\eta^n)/\sqrt D=t^{-n}(1-w^n)/\sqrt D$, so $1/U_n^3=D\sqrt D\,\psi_n$ and \eqref{eq:Xi}
holds.
\end{lemma}
\begin{proof}
(i) For $Q=-1$, $P^2<D<(P+2)^2$ and $D=(P+1)^2$ would give $2P=3$; for $Q=1$, $P\ge3$, $(P-2)^2<D<P^2$ and
$D=(P-1)^2$ would give $2P=5$. (ii) $\sqrt D\ge\sqrt5$ in both cases, and $P\ge1$. (iii) From $\eta^2=P\eta-Q$ we get
$1=Pt-Qt^2$, i.e.\ $t^2-PQt+Q=0$ (as $Q^2=1$); the other root is $Q/t=Q\eta=\sigma t$. (iv) $1/t=P-Qt$, so
$\sqrt D=\eta-\bar\eta=1/t-Qt=P-2Qt$; and $P-2Q^2\eta=P-2\eta=-\sqrt D$. (v) Binet's formula, by induction from
$\eta^2=P\eta-Q$ and $\bar\eta^2=P\bar\eta-Q$; then use $\eta^n=t^{-n}$, $\bar\eta^n=(Qt)^n=w^nt^{-n}$ and
$\sqrt D^3=D\sqrt D$.
\end{proof}

\begin{definition}
For $g$ analytic on a neighbourhood of $|s|\le t^2$ put
$\displaystyle L(g)=\sum_{n\ge1}\rho^n\,[\theta(\theta+1)g](w^n)$.
\end{definition}
The series converges absolutely since $\theta(\theta+1)g=O(s)$ and $|\rho^nw^n|=t^{3n}$.

\begin{lemma}\label{lem:moments}
\textup{(i)} $L(s^k)=\nu_k:=k(k+1)\alpha_k/(1-\alpha_k)$ for $k\ge0$.
\textup{(ii)} For $a\ge1$, $L\bigl(1/(1-w^as)\bigr)=\lambda_a(\Xi-S_a)$ with $\lambda_a=2\rho^{-a}$.
\end{lemma}
\begin{proof}
(i) $\theta(\theta+1)s^k=k(k+1)s^k$; sum the geometric series. (ii) $\theta(\theta+1)(1-cs)^{-1}=2cs/(1-cs)^3$, so the
left side is $\sum_n\rho^n\,2w^{a+n}/(1-w^{a+n})^3=2\rho^{-a}\sum_{N>a}\rho^Nw^N/(1-w^N)^3$, and
$\rho^Nw^N=Q^{2N}t^{3N}=t^{3N}$.
\end{proof}

Let $D_K(s)=\prod_{a=1}^K(1-w^as)$, with poles $x_a=w^{-a}$, $|x_a|=\eta^{2a}$. For a polynomial $P_0$ of degree
$<2K-1$,
\[
\frac{P_0}{D_K}=Q_0+\sum_{a=1}^K\frac{\Res_a(P_0)}{1-w^as},\qquad \Res_a(P_0)=\frac{P_0(x_a)}{\prod_{b\ne a}(1-w^{b-a})},
\]
with $Q_0$ a polynomial of degree $\le\deg P_0-K$. Let $L_X$ be the linear map with $L_X(s^k)=\nu_k$ and
$L_X(1/(1-w^as))=\lambda_a(X-S_a)$, $X$ an indeterminate, and
\[
\Delta_K(X)=\det\bigl[L_X(s^{i+j}/D_K)\bigr]_{0\le i,j<K}\in\QQ(\sqrt D)[X],\qquad\deg\Delta_K\le K.
\]
By Lemma~\ref{lem:moments}, $\Delta_K(\Xi)=\det[L(s^{i+j}/D_K)]$.

\subsection*{The formal parameter}
Every quantity above is the value at $\tau=t$ of a rational function of an indeterminate $\tau$ with rational
coefficients, which depends on $(P,Q)$ only through $Q$: $\rho=Q\tau$, $w=Q\tau^2$, $\alpha_k=Q^{k+1}\tau^{2k+1}$,
$\psi_N=\tau^{3N}/(1-(Q\tau^2)^N)^3$. Write $\Delta_{K,\tau}(X)\in\QQ(\tau)[X]$ for the resulting polynomial, and
$w(\zeta)$, $\alpha_k(\zeta)$ for the values at $\tau=\zeta$. For $Q=-1$ it is the polynomial of \cite{Kri26}. Since
$\sigma t=Q/t$, the coefficientwise conjugate $\sigma\Delta_K$ is the value of $\Delta_{K,\tau}$ at $\tau=Q/t$. Both
values are defined because no denominator vanishes at a real number that is neither $0$ nor a root of unity.

% ===========================================================================
\section{Real decay}\label{sec:decay}
% ===========================================================================

\begin{proposition}\label{prop:decay}
$\log|\Delta_K(\Xi)|\le-\tfrac29(K^3+3K^2-4)\log\eta+O(K^2)$, with an absolute implied constant.
\end{proposition}
\begin{proof}
Let $V_j(s)=\prod_{u=1}^j(s-w^{\lceil u/3\rceil})$, $R_j=\lceil j/3\rceil$ and $a_j=\sum_{u\le j}\lceil u/3\rceil$.
The change of basis from $1,s,\dots,s^{K-1}$ to $V_0,\dots,V_{K-1}$ is unitriangular, so
$\Delta_K(\Xi)=\det[L(V_iV_j/D_K)]$.

Since $\theta(\theta+1)g(s_0)$ depends only on the $2$-jet of $g$ at $s_0$, the $m$-th term of $L(V_iV_jW)$,
$W=1/D_K$, vanishes for $m<R:=\max(R_i,R_j)$: then $V_iV_j$ has a triple zero at $w^m$. For $m\ge R$ and $s=w^m$ every
root $\beta$ of $V_j$ satisfies $|s|\le|\beta|$, so $|s-\beta|\le2|\beta|$ and $\theta(s-\beta)=s$; the signs of the
nodes play no role. Hence $|V_j(s)|\le2^j|w|^{a_j}$ and $|\theta^uV_j(s)|\le j^u2^j|w|^{a_j}|s|/|w|^{R_j}$ for $u=1,2$.
On $|s|\le t^2$ we have $|w^as|\le t^{2a+2}\le\varphi^{-2a-2}$, so $|W|$, $|\theta W|/|s|$ and $|\theta^2W|/|s|$ are
bounded by an absolute constant. Since $\theta(\theta+1)=\theta^2+\theta$ kills constants, by Leibniz every term of
$\theta(\theta+1)(V_iV_jW)$ contains at least one $\theta$ on some factor, hence a factor $\le|s|/|w|^R$. With
$|\rho|^m\le1$ and $\sum_{m\ge R}|w|^{m-R}\le(1-\varphi^{-2})^{-1}$,
\[
|L(V_iV_jW)|\le c_1(i+j+1)^24^{i+j}|w|^{a_i+a_j}
\]
for an absolute constant $c_1$. Pull $|w|^{a_i}$ out of row $i$ and column $i$ and apply Hadamard's inequality:
$|\Delta_K(\Xi)|\le|w|^{2\sum_{j<K}a_j}(c_1(2K)^24^{2K}\sqrt K)^K$. Finally
$2\sum_{j<K}a_j=(K^3+3K^2-\iota_K)/9$ with $\iota_K\in\{0,4,2\}$ according to $K\bmod3$, and $|w|=\eta^{-2}$.
\end{proof}

% ===========================================================================
\section{Denominators}\label{sec:den}
% ===========================================================================

\subsection{The integrality ring and the root types}

\begin{lemma}\label{lem:ring}
The coefficients of $\Delta_{K,\tau}(X)$ lie in $\ZZ[\tau,\tau^{-1},\Phi_d(\tau)^{-1}:d\in\cD_K]$, where $\cD_K$ is
the set of orders of roots of unity $\zeta$ such that
\textup{(W)} $w(\zeta)=Q\zeta^2$ has order $e\le K$, or \textup{($\alpha$)} $\alpha_k(\zeta)=1$ for some $1\le k\le K-2$.
\end{lemma}
\begin{proof}
The inputs are the Laurent monomials $w^a$ and $\lambda_a=2Q^a\tau^{-a}$; the factors $1/(1-w^{b-a})$ and
$(1-w^N)^{-3}$ ($N\le K$), where $1-w^n=1-Q^n\tau^{2n}$ is $\pm$ a product of cyclotomic polynomials $\Phi_d(\tau)$ with
$w(\zeta)^n=1$; and $\nu_k=k(k+1)\alpha_k/(1-\alpha_k)$, of which only $k\le K-2$ occur (the quotient has degree
$\le 2K-2-K$), with $1-\alpha_k=1-Q^{k+1}\tau^{2k+1}=\pm\prod\Phi_d(\tau)$. Division by $D_K$ is integral because its
leading coefficient $\prod_a(-w^a)=\pm\tau^{K(K+1)}$ is a unit. No rational constant is inverted.
\end{proof}

For $d\in\cD_K$ the type depends only on $d$: every primitive $d$-th root is a power $\zeta^j$ with $j$ coprime to $2d$,
and since $Q^j=Q$, $Q\zeta^{2j}=(Q\zeta^2)^j$ and $\alpha_k(\zeta^j)=\alpha_k(\zeta)^j$; powers coprime to the order
preserve the order. If $\zeta$ has order $d$, the order $e(d)$ of $w(\zeta)$ and the least $k\ge1$ with
$\alpha_k(\zeta)=1$ are as follows.
\begin{itemize}
\item[$Q=-1$:] $e(d)=2d$ ($d$ odd), $d$ ($d\equiv2\bmod4$), $d/4$ ($d\equiv4\bmod8$), $d/2$ ($d\equiv0\bmod8$);
$k_{\min}(d)=(d-1)/2$ ($d\equiv3\bmod4$), $(3d-1)/2$ ($d\equiv1\bmod4$), $(d-2)/4$ ($d\equiv2\bmod8$, $d\ge10$),
$(3d-2)/4$ ($d\equiv6\bmod8$), $k_{\min}(2)=2$, and none if $4\mid d$.
\item[$Q=1$:] $e(d)=d$ ($d$ odd), $d/2$ ($d$ even); $\alpha_k(\zeta)=\zeta^{2k+1}=1$ forces $d$ odd, and then
$k_{\min}(d)=(d-1)/2$ for $d\ge3$, $k_{\min}(1)=1$.
\end{itemize}
In both cases $\sum_{e(d)=e}\eu(d)=2\eu(e)$: for each primitive $e$-th root $\omega$ the equation $Q\tau^2=\omega$ has two
roots, and these $2\eu(e)$ roots of unity are exactly those with $e(\tau)=e$. All $d\in\cD_K$ satisfy $d\le4K$:
$w(\zeta)^e=1$ gives $\zeta^{4e}=1$, and $\alpha_k(\zeta)=1$ gives $\zeta^{4k+2}=1$.

\emph{Local setting.} Fix $\zeta$ of order $d\in\cD_K$, let $\varpi=\tau-\zeta$ and let $v=\operatorname{ord}_\varpi$ on
$\QQ(\zeta)((\varpi))$, extended to polynomials in $X$ as the Gauss valuation (the minimum over the coefficients; it is
multiplicative). Then:
\begin{itemize}
\item[(a)] if $e(d)=e\le K$, then $x_a(\zeta)=x_b(\zeta)$ iff $e\mid a-b$, and then $v(x_a-x_b)=1$: indeed
$x_a-x_b=x_a(1-w^{b-a})$, and $1-w^n$ has only simple zeros, since $\frac{d}{d\tau}w^n\ne0$ at $\tau\ne0$;
\item[(b)] $\psi_N$ has a pole of order exactly $3$ if $e\mid N$ and is regular otherwise, by (a); so $S_a$ has a pole of
order $\le3$, and is regular if $a<e$;
\item[(c)] $1-\alpha_k=1-Q^{k+1}\tau^{2k+1}$ is separable, so it has at most a simple zero at $\zeta$;
\item[(d)] $\lambda_a$ and $x_a$ are Laurent monomials, hence units.
\end{itemize}
We call $\zeta$ a \emph{W-root} if $e(d)\le K$ and an \emph{$\alpha$-root} if $\alpha_k(\zeta)=1$ for some
$1\le k\le K-2$; a \emph{pure W-root} is a W-root that is not an $\alpha$-root.

\subsection{Local bounds}

Put $c(m)=2\sum_{j=0}^{m-2}\max(0,\tfrac{m+2}2-j)$ and $c'(m)=2\sum_{j=0}^{m-1}\max(0,\tfrac{m+2}2-j)$; thus
$c(1)=0$, $c(m)=4(m-1)$ for $m\le6$, $c(m)\le(m+3)^2/4$, $c'(1)=3$ and $0\le c'-c\le3$.

\begin{proposition}\label{prop:W}
Let $e=e(d)\le K$ and partition $\{1,\dots,K\}$ into the clusters $I_r=\{a\equiv r\bmod e\}$ of sizes $m_r$,
$1\le r\le e$. Then
\[
-v(\Delta_{K,\tau})\le\sum_{r=1}^{e-1}c(m_r)+c'(m_e)+\delta_\alpha\max(0,K-1-k_{\min}(d)),
\]
with $\delta_\alpha=1$ if $\zeta$ is also an $\alpha$-root and $\delta_\alpha=0$ otherwise.
\end{proposition}
\begin{proof}
The proof of \cite[Proposition~4.2]{Kri26} uses only the facts (a)--(d) and the integrality of the quotient by $D_K$,
which hold for both signs; we repeat it. Let $F(s)=\prod_a(s-x_a)$ and $F_r=\prod_{a\in I_r}(s-x_a)$. Fix a cluster
$r$, write $m=m_r$ and $I_r=\{a_1<\dots<a_m\}$.

\emph{Newton cluster basis.} For $0\le j<m$ put $N_{r,j}=(F/F_r)\prod_{l=m-j+1}^m(s-x_{a_l})$; it vanishes at the top
$j$ nodes of the cluster. At $\varpi=0$ these reduce to $(F/F_r)(s-y_r)^j\cdot$(unit), $y_r=x_r(\zeta)$, which form a
basis of the polynomials of degree $<K$ by the Chinese remainder theorem (distinct residues $r$ give distinct points
$y_r$). So the change of basis from $1,\dots,s^{K-1}$ has unit determinant, and $v(\Delta_{K,\tau})$ is the valuation of
$G=[L_X(NN'/D_K)]$ in this basis.

\emph{Entries.} $G_{(r,j),(r',k)}=P^{\mathrm{pol}}+\sum_a\Res_a(NN')\lambda_a(X-S_a)$, where $P^{\mathrm{pol}}$ is $L_X$ of
the polynomial part. (i) If $r\ne r'$ then $F\mid N_{r,j}N_{r',k}$ and all $\Res_a=0$. (ii) If $r=r'$, then
$\Res_{a_i}=0$ unless $i\le m-\max(j,k)$, and for such $i$
\[
\Res_{a_i}=\frac{(F/F_r)(x_{a_i})^2\prod_{l>m-j}(x_{a_i}-x_{a_l})\prod_{l>m-k}(x_{a_i}-x_{a_l})}
{\prod_{b\ne a_i}(1-w^{b-a_i})}.
\]
Here $(F/F_r)(x_{a_i})$ is a unit, each displayed difference has valuation exactly $1$ by (a), and the denominator has
valuation exactly $m-1$. So $v(\Res_{a_i})=j+k-(m-1)$. (iii) $v(\lambda_a)=0$, and by (b) the Gauss valuation of
$X-S_a$ is $\ge-3$ if $a\ge e$ and $\ge0$ if $a<e$. (iv) For $r<e$ the node $a_1=r$ is regular and $a_i$ is singular
for $i\ge2$; for $r=e$ every node is singular. So the proper part of the $(j,k)$ entry has $v\ge j+k-(m-1)-3$ if some
singular node has $i\le m-\max(j,k)$, and $v\ge j+k-(m-1)$ otherwise. The second case occurs exactly when $r<e$ and
$\max(j,k)=m-1$.

\emph{Weights.} Let $\kappa_j=\max(0,\tfrac{m+2}2-j)$ for $j\le m-2$ and $\kappa_{m-1}=0$ if $r<e$, and
$\kappa_j=\max(0,\tfrac{m+2}2-j)$ for all $j<m$ if $r=e$. Then (A) if the $(j,k)$ entry has a singular node,
$\kappa_j+\kappa_k\ge m+2-j-k$, and (B) if $r<e$ and $\max(j,k)=m-1$, $\kappa_j+\kappa_k\ge0\ge m-1-j-k$. Indeed, in
case (A) if both weights are positive the inequality is an equality; if $\kappa_k=0$, then $k\ge\frac{m+2}2$, so
$\kappa_j\ge\frac{m+2}2-j\ge m+2-j-k$; if both vanish, the right side is $\le0$. Case (B) holds because $j+k\ge m-1$.
Scale row $(r,j)$ by $\varpi^{\lceil\kappa_j\rceil}$ and column $(r,j)$ by $\varpi^{\lfloor\kappa_j\rfloor}$. Since
$a+b\ge n\in\ZZ$ implies $\lceil a\rceil+\lfloor b\rfloor\ge n$, every scaled proper entry is integral. Since
$\kappa_j\in\tfrac12\ZZ$, $\lceil\kappa_j\rceil+\lfloor\kappa_j\rfloor=2\kappa_j$, so the total scaling is
$2\sum\kappa_j$, which is $c(m)$ for $r<e$ and $c'(m)$ for $r=e$. All exponents are $\ge0$.

\emph{Polynomial part.} Write $(\cdot)_+$ for the polynomial part (the quotient by $D_K$). The quotient of $NN'$ has
degree $\le K-2$ and coefficients in $\ZZ[\tau^{\pm1}]$, so $P^{\mathrm{pol}}=\sum_{k\le K-2}[s^k](NN'/D_K)_+\,\nu_k$
(with $\nu_0=0$). At a pure W-root every $\nu_k$, $1\le k\le K-2$, is regular, so the scaled matrix is integral; this
gives the bound with $\delta_\alpha=0$. At an $\alpha$-root, put $P_\alpha=\{1\le k\le K-2:\alpha_k(\zeta)=1\}$ and
write $\nu_k=\varpi^{-1}u_k$ with $u_k$ regular for $k\in P_\alpha$, by (c). Split $P^{\mathrm{pol}}=A^{\mathrm{pol}}+
\varpi^{-1}B$, where $B(f,g)=\sum_{k\in P_\alpha}u_k(0)\,[s^k](fg/D_K)_+$ and $A^{\mathrm{pol}}$ collects
$\varpi^{-1}(u_k-u_k(0))[s^k](\cdot)_+$ for $k\in P_\alpha$ and $\nu_k[s^k](\cdot)_+$ for $k\notin P_\alpha$; both have
regular entries. If $\deg f\le k_{\min}$, then $\deg(fg)\le k_{\min}+K-1$, the quotient of $fg$ has degree
$<k_{\min}$, and $B(f,\cdot)=0$; so $\operatorname{rank}B\le\max(0,K-1-k_{\min})$. After scaling the matrix is
$A+\varpi^{-1}B'$, where $A$ (the scaled proper parts and the scaled $A^{\mathrm{pol}}$) is integral and $B'$ is $B$ in the
Newton basis, scaled; $B'$ has regular entries and $\operatorname{rank}B'=\operatorname{rank}B$, since the basis change
and the scalings are invertible over the fraction field. Expanding the determinant multilinearly in the columns, a term
taking more than $\operatorname{rank}B$ columns from $\varpi^{-1}B'$ vanishes by linear dependence, and every other term
has valuation $\ge-\max(0,K-1-k_{\min})$.
\end{proof}

\begin{lemma}\label{lem:alpha}
If $e(d)>K$ and $\zeta$ is an $\alpha$-root, then $-v(\Delta_{K,\tau})\le\max(0,K-1-k_{\min}(d))$.
\end{lemma}
\begin{proof}
Work in the monomial basis. Since $e>K$, no difference $|a-b|<K$ and no $N\le K$ is divisible by $e$, so no factor
$1-w^{b-a}$ and no $1-w^N$ vanishes at $\zeta$; hence all $\Res_a$, $\lambda_a$ and $S_a$ ($a\le K$) are regular at
$\zeta$, by (c), (d). Only the $\nu_k$, $k\in P_\alpha$, carry poles, and the rank argument of Proposition~\ref{prop:W}
applies with no scaling.
\end{proof}

\begin{proposition}\label{prop:global}
Let $b_{K,d}$ be the bound of Proposition~\ref{prop:W} or Lemma~\ref{lem:alpha} \textup{(}$0$ if $d\notin\cD_K$\textup{)}
and $C_K(\tau)=\prod_d\Phi_d(\tau)^{b_{K,d}}$. Then $C_K(\tau)\Delta_{K,\tau}(X)\in\ZZ[\tau^{\pm1}][X]$; hence
$C_K(t)\Delta_K(X)\in\cO[X]$, and $C_K(\sigma t)\sigma\Delta_K(X)\in\cO[X]$ is its conjugate.
\end{proposition}
\begin{proof}
By Lemma~\ref{lem:ring} a coefficient is $P_1/(\tau^n\prod\Phi_d^{n_d})$ with $P_1\in\ZZ[\tau]$. The bound
$v_\zeta\ge-b_{K,d}$ at every primitive $d$-th root forces $\Phi_d^{n_d-b_{K,d}}\mid P_1$ in $\QQ[\tau]$, hence in
$\ZZ[\tau]$ since $\Phi_d$ is monic. Specialise $\tau\mapsto t$, a unit of $\cO$ ($t=Q\bar\eta=Q(P-\eta)$), and
$\tau\mapsto\sigma t$.
\end{proof}

\subsection{\texorpdfstring{The size of $C_K(t)$}{The size of C\_K(t)}}

\begin{lemma}\label{lem:phid}
For $d\ge1$, $\bigl|\log|\Phi_d(t)\Phi_d(\sigma t)|-\eu(d)\log\eta\bigr|\le2\mathcal E_1$, where
$\mathcal E_1=\sum_{a\ge1}|\log(1-\varphi^{-a})|=2.1136\ldots$.
\end{lemma}
\begin{proof}
$\Phi_d(x)=\prod_{a\mid d}(x^a-1)^{\mu(d/a)}$ and $\sum_{a\mid d}\mu(d/a)a=\eu(d)$. For $|x|=\eta$ write
$x^a-1=x^a(1-x^{-a})$ and use $|\log|1-x^{-a}||\le|\log(1-\eta^{-a})|\le|\log(1-\varphi^{-a})|$; for $x=t$ use
$|1-t^a|$ in the same way.
\end{proof}

Let $\beta_7=\prod_{p\le7}(1-p^{-2})=\frac{768}{1225}$ and $c_*=\frac{11}8\beta_7+\frac13=\frac{4393}{3675}=1.19537\ldots$.

\begin{proposition}\label{prop:cost}
$\log|\Nm C_K(t)|\le\bigl(c_*K^3+O(K^2\log K)\bigr)\log\eta$, with an absolute implied constant.
\end{proposition}
\begin{proof}
By Lemma~\ref{lem:phid}, $\log|\Nm C_K(t)|\le\log\eta\sum_db_{K,d}\eu(d)+2\mathcal E_1\sum_db_{K,d}$, and
$2\mathcal E_1\le(2\mathcal E_1/\log\varphi)\log\eta$. Moreover $\sum_db_{K,d}=O(K^2\log K)$: for each $e\le K$ at most
four orders $d$ have $e(d)=e$ (two if $Q=1$), each with a W part $\le4K+K^2/(4e)+3$ by the cluster bound below, and the
at most $4K$ orders $d\le4K$ have $\alpha$ parts $\le K$; summing $K^2/(4e)$ over $e\le K$ gives $O(K^2\log K)$. Split
$b_{K,d}$ into its W part and its $\alpha$ part.

\emph{W part.} This part depends only on the cluster sizes and on the identity $\sum_{e(d)=e}\eu(d)=2\eu(e)$, which are
the same for both signs, so the estimate of \cite[Proposition~4.6]{Kri26} applies verbatim. Since $m_r\ge1$, $\sum_rm_r=K$,
$m_r<K/e+1$ and $c(m)=4(m-1)+\delta(m)$ with $0\le\delta(m)\le\max(0,m-5)^2/4$,
\[
\sum_{r<e}c(m_r)+c'(m_e)\le4(K-e)+\frac{\max(0,K-4e)^2}{4e}+3.
\]
Bound $\eu(e)\le e\,g_7(e)$ with $g_7(e)=\prod_{p\le7,\,p\mid e}(1-1/p)$, a function of period $210$ with
$0\le g_7\le1$ and mean $\beta_7$. Summing over the orders $d$ with $e(d)=e$ gives the weight
$\sum_{e(d)=e}\eu(d)=2\eu(e)\le2e\,g_7(e)$, so the W part is at most $\sum_{e\le K}g_7(e)f_K(e)$ with
$f_K(y)=8y(K-y)+\tfrac12\max(0,K-4y)^2+6y$. We use the elementary averaging estimate: if $g$ has period $\ell$,
$0\le g\le1$ and mean $\beta$, and $f$ is a function on $[0,K]$ with $|f|\le M$ and total variation $V$, then
\[
\Bigl|\sum_{1\le e\le K}g(e)f(e)-\beta\int_0^Kf(y)\,dy\Bigr|\le2\ell(M+V).
\]
(Cut $[0,K]$ into blocks of length $\ell$. On a full block $(e_0,e_0+\ell]$ both $\sum g(e)f(e)$ and $\beta\int f$
differ from $\ell\beta f(e_0)$ by at most $\ell$ times the oscillation of $f$ on the block; the oscillations add up to at
most $V$, and the last, incomplete block costs at most $2\ell M$.) For $f_K$ we have $M,V=O(K^2)$, so the W part is at
most $\beta_7\int_0^Kf_K+O(K^2)=\tfrac{11}8\beta_7K^3+O(K^2)$.

\emph{$\alpha$ part.} $\max(0,K-1-k_{\min})$ counts the $k\in[k_{\min},K-2]$, so
\[
\begin{aligned}
\sum_d\max(0,K-1-k_{\min}(d))\,\eu(d)&\le\sum_{k=1}^{K-2}\sum_{j=1}^k\#\{\zeta:\alpha_j(\zeta)=1\}\\
&=\sum_{k=1}^{K-2}(k^2+2k)\le\tfrac13K^3+O(K^2),
\end{aligned}
\]
since $\alpha_j(\zeta)=1$ means $\zeta^{2j+1}=Q^{j+1}$, which has $2j+1$ solutions, for either sign.
\end{proof}

\begin{remark}\label{rem:sharp}
The following is a sketch and is not used in the proof. Using $\sum_{n\le y}\eu(n)=3y^2/\pi^2+O(y\log y)$ in residue
classes and the exact $k_{\min}$, the same bound holds with $c_*$ replaced by $(12I+80/27)/\pi^2\le1.130790$ if $Q=-1$
and by $(12I+8/3)/\pi^2\le1.100769$ if $Q=1$. Here $I=\int_0^1\ell(x)x\,dx\in[0.6822834,0.6831236]$, $\ell$ is the
piecewise linear function with $\ell(1/m)=c(m)/m$, and $8/(3\pi^2)$ is the contribution of the odd $d$ with
$k_{\min}=(d-1)/2$.
\end{remark}

% ===========================================================================
\section{Smallness at the conjugate}\label{sec:conj}
% ===========================================================================

\begin{proposition}\label{prop:conj}
$\log|\sigma\Delta_K(-\Xi)|\le-K^2(K+1)\log\eta+K\log(\sqrt K\,M_K)$, where
$M_K=C_0\bigl(K^2/((1-\varphi^{-3})(1-\varphi^{-2}))+\varphi^2\bigr)$ and
$C_0=\prod_{i\ge1}(1-\varphi^{-2i})^{-1}=2.0770906\ldots$.
\end{proposition}
\begin{proof}
At $\tau=\sigma t=Q/t$ the data become: $w\mapsto w^{-1}$, $\lambda_a\mapsto2(\sigma\rho)^{-a}=2t^a$ (as
$\sigma\rho=Q\sigma t=\eta$), $S_a\mapsto-S_a$ (as $\sigma\psi_N=-\psi_N$: $U_N$ is rational and $\sigma\sqrt D=-\sqrt D$),
and $\nu_k\mapsto\sigma\nu_k=Qk(k+1)/(\alpha_k-Q)$ (as $\sigma\alpha_k=Q/\alpha_k$). Let $L^\sigma_X$ be the
corresponding formal functional. The conjugate poles are $r_a=w^a$, $|r_a|=t^{2a}$, and
$\sigma D_K(s)=\prod_a(1-w^{-a}s)=(-1)^Kw^{-A_K}F^\sigma(s)$ with $F^\sigma(s)=\prod_a(s-r_a)$. Hence
\[
\sigma\Delta_K(X)=(-1)^{K^2}w^{A_KK}\det[Y_{i+j}(X)],\qquad Y_k(X)=L^\sigma_X(s^k/F^\sigma).
\]

\emph{The pole moments at $X=-\Xi$.} Since $1/(1-w^{-a}s)=-r_a/(s-r_a)$, $L^\sigma_X(1/(s-r_a))=-r_a^{-1}2t^a(X+S_a)$,
and at $X=-\Xi$ this is $2t^ar_a^{-1}\sum_{N>a}\psi_N=\Gamma(r_a)$, using $t^{4a}=w^{2a}=r_a^2$ (for both signs) and
$w^{a+M}=r_aw^M$, where
\[
\Gamma(r)=2r\sum_{M\ge1}\frac{t^{3M}}{(1-rw^M)^3}=\sum_{n\ge1}\gamma_nr^n,\qquad
\gamma_n=n(n+1)\sum_{M\ge1}t^{3M}w^{M(n-1)},
\]
analytic on $|r|<\eta^2$, with $|\gamma_n|\le n(n+1)t^{2n+1}/(1-t^{2n+1})\le n(n+1)t^{2n+1}/(1-t^3)$. Partial fractions
give
\[
\frac{s^k}{F^\sigma(s)}=\sum_{j=0}^{k-K}h_{k-K-j}(r)\,s^j+\sum_{a=1}^K\frac{r_a^k}{(F^\sigma)'(r_a)(s-r_a)},
\]
and with $\sum_ar_a^m/(F^\sigma)'(r_a)=h_{m-K+1}(r)$ we get
\[
Y_k(-\Xi)=\sum_{j=0}^{k-K}h_{k-K-j}(r)\,\sigma\nu_j+\sum_{n\ge1}\gamma_nh_{k+n-K+1}(r).
\]

\emph{Bounds.} As $h_m(r_1,\dots,r_K)=w^mh_m(1,w,\dots,w^{K-1})$, $|h_m(r)|\le C_0t^{2m}$, using
$\prod_{i\ge1}(1-t^{2i})^{-1}\le C_0$ for $t\le\varphi^{-1}$. For $j\ge1$, $|\sigma\nu_j|\le j(j+1)/(1-t^3)$, so the first
sum is at most $C_0K^2/((1-t^3)(1-t^2))$. For the second,
$\sum_nn(n+1)t^{2n+1}=2t^3/(1-t^2)^3\le2$ and $1/(1-t^3)\le\varphi^2/2$ for $t\le\varphi^{-1}$, so
$\sum_n|\gamma_n|\le\varphi^2$. All these quantities increase with $t$, so $|Y_k(-\Xi)|\le M_K$ for $0\le k\le2K-2$.
Hadamard's inequality and $|w|^{A_KK}=\eta^{-K^2(K+1)}$ finish the proof.
\end{proof}

% ===========================================================================
\section{Nonvanishing}\label{sec:nonvan}
% ===========================================================================

\begin{proposition}\label{prop:nonvan}
$\Delta_K(\Xi)\ne0$ for infinitely many $K$.
\end{proposition}
\begin{proof}
The proof of \cite[Proposition~6.1]{Kri26} uses only $0<|w|<1$, $\rho\ne0$ and the distinctness of the nodes; we repeat
it. Let $\cC(z)=L(1/(z-s))=\sum_n\rho^n\,2w^nz/(z-w^n)^3$ for $|z|>t^2$. It continues to $\CC\setminus\{0\}$ with a pole
of order exactly $3$ at each $w^n$ (leading coefficient $2\rho^nw^{2n}\ne0$), so it is not rational; it is analytic at
$\infty$ and $O(1/z)$. For $\deg R\le n-1$ put $H_{n,R}(z)=L\bigl(R(s)/((z-s)D_n(s))\bigr)$, analytic for $|z|>t^2$
and in particular at every $x_a$, and
\[
T_{n,R}(z)=L\Bigl(\frac{R(z)D_n(s)-D_n(z)R(s)}{(z-s)D_n(s)}\Bigr),
\]
a polynomial in $z$ of degree $\le n-1$. Then $\cC R-T_{n,R}=D_nH_{n,R}$.

If $\Delta_n(\Xi)=0$, any kernel vector gives $R_n\ne0$ of degree $\le n-1$ with $L(s^jR_n/D_n)=0$ for $j<n$. Put
$T_n=T_{n,R_n}$ and $E_n=\cC R_n-T_n$. Expanding $(z-s)^{-1}=\sum_js^jz^{-j-1}$, uniformly with two $s$-derivatives on
$|s|\le t^2$ for $|z|>t^2+\epsilon_0$ ($\epsilon_0>0$), gives $H_{n,R_n}(z)=O(z^{-n-1})$, so $E_n=O(z^{-1})$ and $E_n(x_a)=0$ for $a\le n$.
If $\Delta_n(\Xi)=\Delta_{n+1}(\Xi)=0$, then $V_0=T_{n+1}R_n-T_nR_{n+1}=R_{n+1}E_n-R_nE_{n+1}$ is a polynomial,
$O(z^{n-1})$, vanishing at $x_1,\dots,x_n$; so $V_0=0$ and $T_n/R_n=T_{n+1}/R_{n+1}$. If $\Delta_n(\Xi)=0$ for all $n\ge N$,
all $T_n/R_n$ equal one reduced fraction $T/R_0$, so $R_n=U'_nR_0$, $T_n=U'_nT$, and $E:=\cC R_0-T\ne0$. Then
$U'_nE=E_n=O(1/z)$ forces $E=c_0z^{-q}(1+o(1))$ with $c_0\ne0$, $q\ge1$, and $\deg U'_n\le q-1$; but $E(x_a)\ne0$ for all
large $a$ while $E_n(x_a)=0$, so $U'_n(x_a)=0$ for all large $a\le n$, too many zeros once $n$ is large.
\end{proof}

% ===========================================================================
\section{\texorpdfstring{Proof of Theorem~\ref{thm:main}}{Proof of the theorem}}\label{sec:proof}
% ===========================================================================

Suppose $\sum_nU_n^{-3}=u/m$ with $u\in\ZZ$, $m\in\ZZ_{>0}$. By \eqref{eq:Xi} and Lemma~\ref{lem:unit}(iv),
$\Xi=u\sqrt D/(mD^2)=u(P-2Qt)/(mD^2)$ and $\sigma\Xi=-\Xi$. For $K$ in the infinite set of Proposition~\ref{prop:nonvan}
put $y_K=C_K(t)(mD^2)^K\Delta_K(\Xi)$. By Proposition~\ref{prop:global}, $y_K\in\cO$, and $y_K\ne0$ because
$\Delta_K(\Xi)\ne0$ and $C_K(t)\ne0$ ($t$ is not a root of unity). Moreover
$\sigma y_K=C_K(\sigma t)(mD^2)^K\sigma\Delta_K(-\Xi)$. The norm $y_K\,\sigma y_K$ of the nonzero algebraic integer $y_K$
is a nonzero rational integer, so $|y_K\,\sigma y_K|\ge1$, and by Propositions~\ref{prop:decay}, \ref{prop:conj} and
\ref{prop:cost}
\begin{multline*}
0\le\log|\Nm C_K(t)|+2K\log(mD^2)+\log|\Delta_K(\Xi)|+\log|\sigma\Delta_K(-\Xi)|\\
\le\bigl(c_*-\tfrac29-1\bigr)K^3\log\eta+O(K^2\log K).
\end{multline*}
Since $c_*-\tfrac{11}9=\tfrac{11}8\beta_7-\tfrac89=-0.02684\ldots<0$, this fails for large $K$. \qed

\begin{remark}
Not used in the proof. With the constants of Remark~\ref{rem:sharp} the margin is at least $0.0914$ ($Q=-1$) and
$0.1214$ ($Q=1$) per $K^3\log\eta$. At finite $K$, computed with exact determinants over $\QQ(\sqrt D)$ and interval
arithmetic, the left side above without the term $2K\log(mD^2)$, divided by $K^3\log\eta$, is about $-0.12$ for the Pell
numbers at $K=20$ and $K=25$ and about $-0.14$ for $(P,Q)=(3,1)$, i.e.\ $\sum_nF_{2n}^{-3}$, at $K=20$. With the true
denominator, the norm of the ideal generated by the coefficients of $\Delta_K$, in place of $C_K(t)$, these numbers
become about $-0.17$ and $-0.21$. For $Q=1$ the local bounds of Section~\ref{sec:den} agree with the exact multiplicities
of the cyclotomic factors in $\QQ(\tau)[X]$ for $K\le13$, with equality in most cases. The scripts are available from the
author.
\end{remark}

% ===========================================================================
\section{Formal verification}\label{sec:lean}
% ===========================================================================

% TODO(owner): repository URL once public.
The Lean development (library \lean{LucasU}, Lean~4.35.0, Mathlib v4.35.0, about 9200 lines) generalises the formalisation of
\cite{Kri26} and proves
\begin{center}
\lean{theorem LucasU.irrational\_zetaU (hA : Admissible a E) : Irrational (zetaU a E)},
\end{center}
where \lean{Admissible a E} means $E=1$, $a\ge1$ or $E=-1$, $a\ge3$, and \lean{zetaU a E} is the series
$\sum_{n\ge0}1/U_{n+1}^3$ for the sequence \lean{Useq a E} with $U_{n+2}=aU_{n+1}+EU_n$, written with a \lean{tsum}; so
$(a,E)=(P,-Q)$. The proof uses only the axioms \lean{propext}, \lean{Classical.choice} and \lean{Quot.sound}; there is no
\lean{sorry}, no \lean{native\_decide} and no added axiom.

\emph{What has to be trusted.} For the theorem itself, only the definitions of \lean{Useq}, \lean{zetaU} and
\lean{Admissible} (and Mathlib's \lean{tsum} and \lean{Irrational}) have to be read. Since a divergent \lean{tsum} has
the value $0$, which is rational, the statement cannot hold for a degenerate reason; moreover \lean{hasSum\_zetaU}
shows that the series converges (to $D\sqrt D\,\Xi$). The corollaries are \lean{LucasU.irrational\_pell}, that the
\lean{tsum} of $1/\cP_{n+1}^3$ over $n\ge0$ is irrational, with \lean{pell} $=$ \lean{Useq 2 1}, and
\lean{LucasU.irrational\_fib\_mul}, that for $k\ge1$ the \lean{tsum} of $1/F_{k(n+1)}^3$ over $n\ge0$, written with
\lean{Nat.fib}, is irrational; the identity $F_{kn}=F_kU_n(L_k,(-1)^k)$ is \lean{LucasU.fib\_mul\_eq}, stated as
$F_{kn}=F_k\,U_n$ for \lean{Useq (lucasNum k) ((-1)\^{}(k+1))}, i.e.\ $(a,E)=(L_k,(-1)^{k+1})$, with the
Lucas numbers \lean{lucasNum} (all in \lean{LucasU/Corollaries.lean} and \lean{LucasU/Final.lean}).

\subsection*{Definitions and the chain}
\lean{LucasU/Defs.lean} defines the construction of Section~\ref{sec:setup} over an arbitrary field with a sign
$\varepsilon=Q$ and a parameter $\tau$ (\lean{wpar} $=\varepsilon\tau^2$, \lean{alpha}, \lean{nu}, \lean{psi}, \lean{S},
\lean{node}, \lean{F}, \lean{res}, \lean{mConst}, \lean{mLin}, \lean{nrm}, \lean{Delta}), the real functional \lean{Lfun}
and \lean{XiR} as functions of $(\varepsilon,t)$, and the clearing polynomial \lean{CKpoly} with exponents \lean{nExp}.
\lean{LucasU/Statements.lean} states the ingredients as propositions in $(\varepsilon,t)$ (and \lean{Integrality} in
$\varepsilon$ alone, for every real $\tau\ne0$ that is not a root of unity); the proofs below assume
\lean{Good}~$\varepsilon$~$t$, i.e.\ $\varepsilon=\pm1$ and $0<t\le\varphi^{-1}$. The conjugate point is
$\tau=\varepsilon/t$. \lean{LucasU/Unit.lean} contains
Lemma~\ref{lem:unit}, and \lean{LucasU.irrational\_zetaU\_of} (in \lean{LucasU/Main.lean}) derives the theorem from the
ingredients at $\varepsilon=-E$, $t=1/\eta$, using that $t$ and $\varepsilon/t$ are the two roots of
$\tau^2+aE\tau-E$ and that $\sqrt D$ is irrational (\lean{norm\_ne\_zero}).
\begin{center}
\small
\begin{tabular}{lll}
\toprule
Paper & Lean statement & Lean proof\\
\midrule
Lemma~\ref{lem:moments}, entries & \lean{EntryIdentity} & \lean{LucasU.Analytic.entryIdentity}\\
Prop.~\ref{prop:decay} & \lean{DecayBound} & \lean{LucasU.Analytic.decayBound}\\
Lemma~\ref{lem:ring}--Prop.~\ref{prop:global} & \lean{Integrality} & \lean{LucasU.Denom.integrality}\\
Lemma~\ref{lem:phid}, Prop.~\ref{prop:cost} & \lean{CostBound} & \lean{LucasU.Cost.costBound}\\
Prop.~\ref{prop:conj} & \lean{ConjBound} & \lean{LucasU.Conj.conjBound}\\
Prop.~\ref{prop:nonvan} & \lean{Nonvanishing} & \lean{LucasU.NonVan.nonvanishing}\\
Section~\ref{sec:proof} & & \lean{LucasU.irrational\_zetaU\_of}\\
\bottomrule
\end{tabular}
\end{center}
\lean{EntryIdentity} is used only inside the proof of \lean{DecayBound}, not by the assembly. The formal statements are
the propositions in the form the assembly needs, with constants that may depend on $(\varepsilon,t)$: the uniformity
in $(P,Q)$ claimed in Propositions~\ref{prop:decay}, \ref{prop:cost} and the explicit $M_K$ of
Proposition~\ref{prop:conj} are not expressed by the formal statements and are not needed, since the theorem
concerns one fixed $(P,Q)$ (the formal proof of \lean{DecayBound} does give the uniform constant $C=4096$).
\lean{DecayBound} is $|\Delta_K(\Xi)|\le C^{K^2}t^{2b_K}$ with $b_K=2\sum_{j<K}a_j$ (proved with $C=4096$),
\lean{ConjBound} is $|\sigma\Delta_K(-\Xi)|\le C^{K^2}t^{K^2(K+1)}$, and \lean{CostBound} is
$\log|C_K(t)C_K(\varepsilon/t)|\le\frac{121}{100}K^3\log t^{-1}+CK^2(1+\log K)^2$; its proof follows
Proposition~\ref{prop:cost} with the cubic constant $c_*=\frac{4393}{3675}$ for both signs, and
$c_*<\frac{121}{100}<\frac{11}9$. The proof of \lean{Integrality} works in the local rings of $\CC(\tau)$ at the roots
of unity, and uses the root types only through \lean{LucasU/Roots.lean}; the proof of \lean{Nonvanishing} works with the
Cauchy transform $\cC$ directly.

\subsection*{Provenance}
The formal development generalises that of \cite{Kri26}, parts of which are adapted from the Lean formalisation QZeta3
\cite{QZ3} of \cite{UPL} (Apache License~2.0); every adapted file says so in its header.

\subsection*{Use of AI tools}
This work was carried out with substantial assistance from AI systems (Anthropic's Claude models), which were used to
search the literature, to develop and check the proof, to write the Lean formalisation and to draft this text. The
author directed the project and is responsible for its content. The formal proof is checked by the Lean kernel.

% ===========================================================================
\bibliographystyle{amsplain}
\bibliography{refs}

\end{document}
